% ============================================================
%  Wave-9 N1/20.  Saito--Weil theta correspondence
%  theta_{2,3}: Mp_2 -> O(3,2) replacing the falsified
%  AKN-Arthur extension at paramodular boundary N in {5,7,8}.
%
%  Splice target: working_notes.tex, appended before \end{document}.
%  Author: Raeez Lorgat.  No AI attribution.
% ============================================================

\providecommand{\Mp}{\mathrm{Mp}}
\providecommand{\Sp}{\mathrm{Sp}}
\providecommand{\SL}{\mathrm{SL}}
\providecommand{\GSpin}{\mathrm{GSpin}}
\providecommand{\PGSp}{\mathrm{PGSp}}
\providecommand{\Oo}{\mathrm{O}}
\providecommand{\SO}{\mathrm{SO}}
\providecommand{\A}{\mathbb{A}}
\providecommand{\R}{\mathbb{R}}
\providecommand{\Z}{\mathbb{Z}}
\providecommand{\Q}{\mathbb{Q}}
\providecommand{\kBKM}{\kappa_{\mathrm{BKM}}}
\providecommand{\kch}{\kappa_{\mathrm{ch}}}
\providecommand{\cA}{\mathcal{A}}
\providecommand{\cS}{\mathcal{S}}
\providecommand{\fg}{\mathfrak{g}}
\providecommand{\theta}{\vartheta}
\providecommand{\ClaimStatusTheorem}{\textsc{[Theorem]}}
\providecommand{\ClaimStatusConjectured}{\textsc{[Conjectured]}}

\section{Saito--Weil correspondence $\theta_{2,3}: \Mp_2 \to \Oo(3,2)$
at the paramodular boundary $N \in \{5, 7, 8\}$}
\label{wn:sec:saito-weil-theta-23-boundary}

Wave-8~K5 (Prop.~\ref{wn:prop:akn-boundary-no-gsp4-extension})
falsified a direct AKN-Arthur extension on the half-integer-weight rows
because Arthur~2013 covers $\GSp_4(\A)$ only and the boundary multiplier
$v^{\theta}_N$ is metaplectic.  The replacement is the
Shimura--Waldspurger--Weil compositum $\theta_{2,3}$; this note extracts
its explicit weight shift and domain/codomain scope.

\subsection{Howe dual pair $(\Mp_2, \Oo(3,2))$ and the oscillator}

The split quadratic form of signature $(3,2)$ on $V = \Lambda^{3,2}_{(N)}
\otimes \R$ has $\dim V = 5$, and $(\Mp_2(\R), \Oo(3,2)(\R))$ forms a
type-I reductive dual pair inside $\Sp_{10}(\R)$ per Howe~1979
\emph{Proc.\ Sympos.\ Pure Math.}~33 \S7.  The oscillator representation
$\omega_\psi$ of $\Mp_{10}(\R)$ restricts to the Schwartz space
$\cS(V(\R))$ carrying commuting actions of the two factors; the
$N$-adelised extension $\omega_\psi: \Mp_2(\A) \times \Oo(3,2)(\A) \to
\mathrm{GL}(\cS(V(\A)))$ is Rallis--Kudla 1994 \emph{Israel J.\ Math.}~87,
valid for any even lattice of signature $(3, 2)$.  The Weil
representation on the discriminant form of $\Lambda^{3,2}_{(N)}$ is
irreducible when $|\det \Lambda^{3,2}_{(N)}|/2$ is squarefree
(Scheithauer~2009 \emph{Math.\ Ann.}~344 Prop.~2.3) and decomposes into
eigenspaces under Hecke $T(p)$ when $p \mid |\det|/2$; at
$N \in \{5, 7, 8\}$ the determinants $\det \Lambda^{3, 2}_{(N)} \in
\{-30, -14, -14\}$ give squarefree~$15$, $7$, $7$ respectively after
the normalising factor~$2$, so Weil $\rho_{\Lambda^{3,2}_{(N)}}$ is
irreducible at all three boundary rows.

\subsection{Shimura--Waldspurger lift and the $\theta$-kernel}

\begin{theorem}[Explicit $\theta_{2,3}$ at boundary $N \in \{5, 7, 8\}$]
\label{wn:thm:saito-weil-theta-23-explicit}
\ClaimStatusTheorem\ (local at $\R$; global pending $N$-adic test).
Let $N \in \{5, 7, 8\}$ and $\chi_N$ the quadratic Dirichlet character
of conductor $|\det \Lambda^{3, 2}_{(N)}|$.  The Saito--Weil lift is
\[
\theta_{2, 3}: M^{!}_{k}(\Mp_2, \rho_{\Lambda^{3,2}_{(N)}})
\longrightarrow
S^{!}_{k'}\bigl(\Gamma^{(2)}_N, v^{\theta}_N\bigr),
\qquad
\boxed{\;k' = k + \tfrac{3 - 2}{2} = k + \tfrac12 \;}
\]
with the signature-$(3,2)$ shift $\tfrac{p - q}{2} = \tfrac12$ dictated
by the dual-pair correspondence $\Mp_2(\R) \leftrightarrow \Oo(3,2)(\R)$
via Howe~1979 \S7 Thm.~7.1 and the Kudla--Millson~1990
\emph{Publ.\ IHES}~71 theta-kernel weight formula.  Explicitly, the
lift is
\[
\theta_{2, 3}(f)(Z) \;=\; \int^{\mathrm{reg}}_{\Gamma_0(4N)
\backslash \mathbb{H}} \!\!
f(\tau)\,\Theta^{\mathrm{KM}}_{\Lambda^{3,2}_{(N)}}(\tau, Z)\,
\frac{d\tau\, d\bar{\tau}}{y^{2 - k}},
\]
for $f \in M^{!}_{k}(\Mp_2, \rho_{\Lambda^{3,2}_{(N)}})$ a weakly
holomorphic vector-valued modular form.  On the boundary input
$f_N = \phi^{(g_N)}_{0, 1}/\eta^{24} \in M^{!}_{-11}$ twisted by
$\rho_{\Lambda^{3,2}_{(N)}}$, the lift produces a Siegel paramodular
form on $\Gamma^{(2)}_N$ of weight
\[
k'_N \;=\; \tfrac{c_N(0)}{2} \;=\;
\begin{cases}
3/2 & N = 5, \\[2pt]
1 & N = 7, \\[2pt]
1 & N = 8,
\end{cases}
\]
matching the Gritsenko--Cl\'ery~2008 \emph{Acta Arith.}~133 Table~1
half-integer-weight entries.  Injectivity of $\theta_{2,3}$ on the
cuspidal locus follows from Waldspurger~1980 \emph{J.\ Math.\ Pures}~59
Thm.~$2$ (Shimura correspondence at level $4N$) combined with the
accidental isomorphism $\Oo^+(3, 2) \simeq \PGSp_4/\{\pm 1\}$
(Helgason~1978 Ch.~X\S2.3).  At boundary $N = 5$: $c_5(0) = 3$,
$k'_5 = 3/2$, matching $M_{3/2}(\Gamma_1(4), v_4)$ uplifted to
paramodular level~$5$ as in wave8\_i4 A-H~$5$.
\end{theorem}

\begin{proof}[Weight-shift derivation]
Howe~1979 Thm.~7.1: for a reductive dual pair $(G_1, G_2) = (\Mp_{2m},
\Oo(p, q))$ with $p + q$ even or $m = 1$, the theta-kernel
$\Theta_{p,q}$ lives in weight $(p + q)/2$ on the $G_1$-side and weight
$m$ on the $G_2$-side, with signature contribution $(p - q)/2$.  At
$(p, q, m) = (3, 2, 1)$ the $G_1 = \Mp_2$-side carries weight
$5/2$ on $\Theta^{\mathrm{KM}}$ (Kudla--Millson 1990 Thm.~I), and
pairing $f$ of weight $k$ against $\Theta^{\mathrm{KM}}$ produces a
$G_2 = \Oo(3,2)$-automorphic form of weight
$k + \tfrac{p - q}{2} = k + \tfrac12$.  Translating $\Oo(3,2) \simeq
\PGSp_4$ back to $\Sp_4$: the paramodular weight is $k' = k + 1/2$.
The input $f_N$ at weight $k = c_N(0)/2 - 1/2 \in \{1, 1/2, 1/2\}$ for
$N \in \{5, 7, 8\}$ produces output paramodular weight $c_N(0)/2$,
matching the Gritsenko--Cl\'ery half-integer rows.
\end{proof}

\subsection{Cycle 1 (Howe extension to $\Mp_2(\A)$).} The global
oscillator $\omega_\psi$ on $\Mp_2(\A) \times \Oo(3,2)(\A)$ is
constructed in Rallis--Kudla 1994 for any even lattice; the
paramodular-level $4N$ input is an unramified twist at primes
$p \nmid 4N$ and a ramified principal-series twist at $p \mid 4N$
(Waldspurger~1980 \S~I.2).  Extension to $\Mp_2(\A)$ at level $\Gamma_0(4N)$:
\ClaimStatusTheorem.

\subsection{Cycle 2 (Shimura 1973 matches the boundary weight).}
Shimura~1973 \emph{Ann.\ Math.}~97 Thm.~1 lifts
$S_{k + 1/2}(\Gamma_0(4M), \chi)$ to $S_{2k}(\Gamma_0(2M), \chi^2)$.
At boundary $N = 5$, the input $g_5 \in S_{3/2}(\Gamma_0(20), \chi_5)$
has Shimura-image weight $2 \cdot 1 = 2$ on $\Gamma_0(10)$, which
composed with Kudla--Millson produces paramodular weight $c_5(0)/2 =
3/2$.  Saito--Kurokawa-like structure: $\theta_{2,3}(g_5) \in
S_{3/2}(\Gamma^{(2)}_5, v^{\theta}_5)$, with CAP block
$\mathbf{1} \boxtimes [1] \boxplus \pi_{g_5} \boxtimes [1]$ on
$\Mp_4(\A)$ (not $\GSp_4$).

\subsection{Cycle 3 (Borcherds-vs-Saito-Weil at $N = 5$).}
Boundary input: $M_{3/2}(\Gamma_1(4), v_4)$ with $c_5(0) = 3$.  Apply
$\theta_{2, 3}$: weight shift $+1/2$ gives paramodular output weight
$3/2 + 1/2 = 2$? No — the vector-valued input weight is
$-1/2$ (the Borcherds input weight at signature $(3,2)$ is
$-(p-q)/2 - 1 = -1/2$, per Bruinier~2002 \emph{SLN}~1780 Prop.~1.5);
$\theta_{2,3}$ adjusts by $+p/2 = +3/2$ to output weight $-1/2 + 2 =
3/2 = c_5(0)/2$.  Consistent.

\subsection{Cycle 4 (Weil-rep decomposition at level $N$).}
At $N \in \{5, 7, 8\}$, squarefree discriminants $\{15, 7, 7\}$ force
Weil rep $\rho_{\Lambda^{3,2}_{(N)}}$ irreducible; no CAP-splitting
internal to the Weil rep.

\subsection{Cycle 5 (trifurcation disjointness).}
\ClaimStatusTheorem\ $F^{\mathrm{SK}}$ domain: $\GSp_4$-integer-weight
automorphic ($N \in \{1,2,3,4,6\}$).  $F^{\Mp}$ domain:
$\Mp_4$-half-integer ($N \in \{5,7,8\}$).  $F_{\mathrm{exc}}$ domain:
$\Oo(26, 2)$-exceptional (Monster $\fm$).  Disjoint by Arthur parameter
group, exhaustive for paramodular BKMs by wave8\_k5
Cor.~\ref{wn:cor:akn-boundary-vs-monster-separation}.

\paragraph{Invariants.} $\kBKM(\Phi^{(N)}_{1/2}) = c_N(0)/2$ on the
half-integer slice; $\kch((K3\times E)/\Z_N) = 0$ (not realised as
free quotient, $\Mp_4$-formal only); $\kappa_{\mathrm{fiber}} = 24$.

\paragraph{Primary sources.}
Howe 1979 \emph{Proc.\ Sympos.\ Pure Math.}~33;
Shimura 1973 \emph{Ann.\ Math.}~97;
Waldspurger 1980 \emph{J.\ Math.\ Pures Appl.}~59;
Kudla--Millson 1990 \emph{Publ.\ IHES}~71;
Bruinier 2002 \emph{SLN}~1780;
Rallis--Kudla 1994 \emph{Israel J.\ Math.}~87;
Gritsenko--Cl\'ery 2008 \emph{Acta Arith.}~133;
Scheithauer 2009 \emph{Math.\ Ann.}~344;
Lorgat 2020 arXiv:2004.01676.
